\begin{proof}\textbf{of Theorem~\ref{thm: PTF MSE lower bound}.}
For each $i$, we can still define the minimum norm interpolant as
\[
g_i
=
\argmin_{g\in\mathcal H_K}
\left\{
\|g\|_{\mathcal H_K}:
g(t_j)=f_i(t_j),\ j=1,\ldots,T
\right\},
\]
and its transformed coefficients as in the proof for Theorem~\ref{thm: FTP MSE RKHS}:
\[
\mathbf \beta_i
:=
K_T^{-1/2}\mathbf f_i,\quad
\|g_i\|_{\mathcal H_K} = \|\beta_i\|_2\le B.
\]
Denote $a_j = \sqrt{K(t_j,t_j)}$, by the RKHS reproducing property, 
\[
|f_i(t_j)|
\leq
\|f_i\|_{\mathcal{H}_K}\sqrt{K(t_j,t_j)}
\leq
Ba_j
\leq 
B' a_j.
\]
Since we are considering the noiseless case $\sigma_0=0$ and $x_i = \mathbf{f}_i$, with
\[
c_j = B'a_j.
\]
the clipping operation $\operatorname{clip}(x_{ij},c_j)$ is inactive and
\[
\bar x^{(\mathrm{clip})}
=
\bar{x}
=
\dfrac{1}{n}\sum_{i=1}^n x_i
=
\bar{\mathbf{f}}.
\]
It follows that the released private vector can be written as
\[
\widetilde x
=
\bar{\mathbf{f}}+Z
=
K_T^{1/2}\bar{\beta} + Z
,
\]
where
\[
Z\sim N(0,\tau^2W^{-1}),
\quad
\tau^2=\frac{4B'^2T}{n^2\mu^2},
\quad
W = \mathrm{diag}\{a_1^{-2},\ldots,a_T^{-2}\}.
\]
By the representer theorem, the constrained estimator has the form
\[
\hat f_{\mathrm{PTF}}^{(c)}(t)
=
\sum_{j=1}^T
(\hat\alpha^{(c)})_jK(t_j,t),
\]
\[
\hat\alpha_{\mathrm{PTF}}^{(c)}
=
\argmin_{\alpha^\top K_T\alpha\le {B'}^2}
\dfrac{1}{T} 
(\widetilde{x} - K_T\alpha)^\top W (\widetilde{x} - K_T\alpha).
\]
Denote $\beta = K_T^{1/2}\alpha$, we can rewrite the constrained problem as
\[
\hat\beta_{\mathrm{PTF}}^{(c)}
=
\argmin_{\|\beta\|_2\le B'}
(\widetilde{x} - K_T^{1/2}\beta)^\top W (\widetilde{x} - K_T^{1/2}\beta),
\quad 
\hat\beta_{\mathrm{PTF}}^{(c)} = K_T^{1/2}\hat\alpha_{\mathrm{PTF}}^{(c)}.
\]
We then analyze the corresponding unconstrained WLS estimator.\\
Define 
\[
\Phi:=K_T^{1/2},
\quad
M:=\Phi^\top\Phi = K_T,
\quad
A:= \Phi^\top W\Phi
=K_T^{1/2}WK_T^{1/2}.
\]
By
definition, the unconstrained problem is:
\[
\hat{\beta}_{\mathrm{PTF}}^{(u)}
=
\argmin_{\beta\in \R ^T} (\widetilde{x} - K_T^{1/2}\beta)^\top W (\widetilde{x} - K_T^{1/2}\beta).
\]
with the solution
\[
\widetilde\beta_{\mathrm{PTF}}^{(u)}
=
A^{-1}\Phi^\top W\widetilde x
=
\Phi^{-1}\widetilde{x}.
\]
Substituting 
$\widetilde x=K_T^{1/2}\bar{\beta} + Z= \Phi\bar\beta + Z$ gives $\widetilde\beta_{\mathrm{PTF}}^{(u)}=\bar\beta+\xi$, where $\xi:=\Phi^{-1}Z$.\\
Since $Z\sim \mathcal{N}(0,\tau^2W^{-1})$, we obtain
\[
\xi
\sim
\mathcal{N} (0,\tau^2 \Phi^{-1}W^{-1}\Phi^{-1})
\overset{d}{=}
\mathcal{N}(0,\tau^2 A^{-1}).
\]
Hence
\begin{equation}
\mathbb E
\left\|
\xi
\right\|_2^2
=
\tau^2\operatorname{tr}(A^{-1}),\quad
\mathbb E
\left\|
Z
\right\|_2^2
=
\tau^2\operatorname{tr}(W^{-1}).
\label{eq: PTF unconstrained MSE_RKHS}
\end{equation}
We next account for the constraint $\|\beta\|_2\le B'$. 
Let $\mathcal A=\{\|\widetilde\beta_{\mathrm{PTF}}^{(u)}\|_2>B'\}$ denote the event that the unconstrained estimator violates the constraint. On the complementary event $\mathcal A^c$, the unconstrained estimator is feasible and hence
$\widetilde\beta_{\mathrm{PTF}}^{(c)}
=
\widetilde\beta_{\mathrm{PTF}}^{(u)}$. Therefore
\[
\begin{aligned}
\mathbb E\|\Phi(\widetilde\beta_{\mathrm{PTF}}^{(c)}-\bar\beta)\|_2^2
&=
\mathbb E[
\|\Phi(\widetilde\beta_{\mathrm{PTF}}^{(u)}-\bar\beta)\|_2^2
\mathbf 1_{\mathcal A^c}]
+
\mathbb E[
\|\Phi(\widetilde\beta_{\mathrm{PTF}}^{(c)}-\bar\beta)\|_2^2
\mathbf 1_{\mathcal A}]
\\
&\geq
\mathbb E[
\|\Phi\xi\|_2^2
\mathbf 1_{\mathcal A^c}]
\\
&=
\mathbb E\|Z\|_2^2
-
\mathbb E[
\|Z\|_2^2
\mathbf 1_{\mathcal A}].
\end{aligned}
\]
Thus, by Equation~\eqref{eq: PTF unconstrained MSE_RKHS},
\begin{equation}
\begin{aligned}
    \operatorname{MSE}_{\mathrm{PTF}}
    =
    \dfrac{1}{T}\mathbb{E} 
    \|\widetilde{\mathbf{f}}_{\mathrm{PTF}} - \bar{\mathbf{f}}\|_2^2
    =
    \dfrac{1}{T}\mathbb{E} 
    \|\Phi(\widetilde\beta_{\mathrm{PTF}}^{(c)}-\bar\beta)\|_2^2
    \geq
    \dfrac{1}{T}\SqBra{\tau^2\operatorname{tr}(W^{-1})
    -
    R_{\mathrm{PTF}}},
\end{aligned}
\label{eq: PTF remainder decomposition_RKHS}
\end{equation}
where
\[
R_{\mathrm{PTF}}
=
\mathbb E[
\|Z\|_2^2
\mathbf 1_{\mathcal A}].
\]
It remains to bound the remainder term. Since
\[
\|\bar\beta\|_2\le \dfrac{1}{n}\sum_{i=1}^n\|\beta_i\|_2\leq B,
\]
the event $\mathcal A$ implies
\[
\|\xi\|_2
=
\|\widetilde\beta_{\mathrm{PTF}}^{(u)}-\bar\beta\|_2
\geq
\|\widetilde\beta_{\mathrm{PTF}}^{(u)}\|_2-\|\bar\beta\|_2
>
B'-B.
\]
Consequently,
\[
\mathcal A
\subseteq
\left\{
\|\xi\|_2>B'-B
\right\}.
\]
Write $\xi=\tau A^{-1/2}z,\ z\sim N(0,I_T)$. Since
\[
\|A^{-1/2}z\|_2^2
\leq
\lambda_{\max}(A^{-1})\|z\|_2^2
=
\frac{\|z\|_2^2}
{\lambda_{\min}(A)},
\]
we obtain
\[
\Pr(\mathcal A)
\leq
\Pr\left(
\|z\|_2^2
>
\frac{(B'-B)^2\lambda_{\min}(A)}
{\tau^2}
\right).
\]
Because $\|z\|_2^2\sim\chi_T^2$, apply the chi-square tail bound in Equation~\eqref{eq: chi_square bound}, we get
\begin{equation}
\Pr(\mathcal A)
\leq
\exp\left[
-\frac12
\left\{
\frac{(B'-B)\sqrt{\lambda_{\min}(A)}}{\tau}
-\sqrt T
\right\}^2
\right].
\label{eq: PTF active probability_RKHS}
\end{equation}
Finally, by the Cauchy--Schwarz inequality,
\[
R_{\mathrm{PTF}}
\leq
\left(
\mathbb E\|Z\|_2^4
\right)^{1/2}
\Pr(\mathcal A)^{1/2}.
\]
For a centered Gaussian vector with covariance
$\tau^2W^{-1}$,
\[
\mathbb E\|Z\|_2^4
=
\tau^4
\left[
\{\operatorname{tr}(W^{-1})\}^2
+
2\operatorname{tr}(W^{-2})
\right]
\le 
\tau^4 [3\{\operatorname{tr}(W^{-1})\}^2].
\]
Therefore
\[
\left(
\mathbb E\|Z\|_2^4
\right)^{1/2}
\leq
\tau^2
\sqrt{3}\operatorname{tr}(W^{-1}).
\]
Combining this inequality with Equation~\eqref{eq: PTF active probability_RKHS}
and Equation~\eqref{eq: PTF remainder decomposition_RKHS}, we prove 
Equation~\eqref{eq: PTF MSE lower bound}.
\end{proof}

\begin{proof}
\textbf{of Corollary~\ref{col: MSE ratio in RKHS}.}
Write $\tau_n = \tau =2B_n'\sqrt T/(n\mu)$. Since $B_n'$ is a non-decreasing sequence  satisfying $B_n'>B$ for sufficiently large $n$, we have
\[
\frac{B_n'-B}{\tau_n}
=
\frac{n\mu}{2\sqrt T}
\left(1-\frac{B}{B_n'}\right)
\longrightarrow\infty.
\]
Thus, the radius condition in Theorem~\ref{thm: PTF MSE lower bound} holds for all
sufficiently large $n$, and its exponential remainder tends to zero. Combining that lower bound with the exact noiseless FTP error in Theorem~\ref{thm: FTP MSE RKHS} gives
\[
\limsup_{n\to\infty}
\frac{\operatorname{MSE}(\widetilde f_{\mathrm{FTP}})}
     {\operatorname{MSE}(\widetilde f_{\mathrm{PTF}})}
\le
\frac{\operatorname{tr}(K_T)}
     {T\operatorname{tr}(W^{-1})}.
\]
By the definition
$W=\operatorname{diag}\{K(t_j,t_j)^{-1}:j=1,\ldots,T\}$,
\[
\operatorname{tr}(W^{-1})
=\sum_{j=1}^T K(t_j,t_j)
=\operatorname{tr}(K_T).
\]
Hence, we get the claimed bound of $1/T$.
\end{proof}